%
\againframe{plan}
%
\begin{frame}{Lecture plan}
	\begin{learningobjectives}
		\begin{itemize}
			\item From qualitative picture of metals to quantitative derivation.
			\item Derive expression for conductivity.
			\item Mobility is a material property.
		\end{itemize}
	\end{learningobjectives}
\end{frame}
%
%
\begin{frame}{Free electron density -- \(n\)}
	\begin{block}{Definition}
		\begin{itemize}
			\item Electron density \(n\) is defined as the number of electrons per unit volume. Its SI unit is \si{\per\metre\cubed}.
			\item For metals, each atom contributes as many electrons as its \textbf{valency}. For example, \ce{Na} has valency 1.
		\end{itemize}
	\end{block}
	\begin{guesstimate}{Electron density of metals}
		%
		\begin{enumerate}\itemsep0pt\parskip0pt \parsep0pt
			\item Take the radius \(r_s\) of atom is nearly $3 a_0$, where \(a_0 =\)~\SI{0.529}{\angstrom} is the Bohr radius.
			\item Deduce \(n\) is of the order \SI{e22}{\per\centi\metre\cubed} =  \SI{e28}{\per\metre\cubed}.
			\item For Cu, the experimental observed \(n\)~=\SI{8.47e28}{\per\metre\cubed}.
		\end{enumerate}
	\end{guesstimate}
	%
	\vfill
	\begin{keyinsight}
		\(\tfrac{r_s}{a_0}\) determines the free electron density.
	\end{keyinsight}
	%
\end{frame}
%
\begin{frame}{Drift velocity -- \(v_d\)}
	\begin{minipage}{0.7\textwidth}
		\begin{definition}
			\begin{itemize}
				\item \textbf{Drift velocity} \(\vb{v_d}\)is defined as the average velocity of an electron in the presence of electric field.
				\item \textbf{relaxation time} $\tau$ is defined as the average time between collisions.
			\end{itemize}
		\end{definition}
	\end{minipage}
	\begin{minipage}{0.28\textwidth}
		\begin{figure}
			\centering
			\includegraphics[height=0.05\textheight]{metals_drift_equation}
		\end{figure}
	\end{minipage}
	\begin{derivation}
		\begin{minipage}[t]{0.6\textwidth}
			\begin{itemize}
				\item The equation of motion of electron is
				\item During time duration \(\tau\), electron reaches drift velocity
			\end{itemize}
		\end{minipage}
		%
		\begin{minipage}[t]{0.38\textwidth}
			\[
			\vb{a} = -\frac{e}{m} \vb{E}
			\]
			\[		
			\framebox{
				$
				\vb{v_d} = -\frac{e\tau}{m}\vb{E}
				$}
			\]
		\end{minipage}
	\end{derivation}
	\begin{keyinsight}
		Drift velocity is proportional to the applied electric field.
	\end{keyinsight}
\end{frame}
%
\begin{frame}{\(y=mx\) and \(y= m /x\) relations}
	\begin{itemize}
		\item We have seen many examples of \textbf{linear} relations between physical quantities till now.
		\item The dependent variable \(y\) relates to the independent variable \(x\). The relation is given by the \textbf{constant of proportionality}.
		\item We also say that ``\(y\) is directly proportional to \(x\)''.
	\end{itemize}
	%
{
	%\small
		\begin{table}
		\centering
		\begin{tabular}{l c l}
			\toprule
			Law 	&	Equation 	&	Const. of prop.	\\
			\midrule
			Ohm's (macroscopic)	&	\(I = \frac{1}{R} V\)	&	Conductance	\\
			Ohm's (microscopic)	&	\(\vb{j} = \sigma \vb{E}\)	&	Conductivity \\
			Capacitor			&	\(E = \frac{1}{d} V\)	&	1/Separation	\\
			Drift motion			&	\(\vb{v_d} = -\frac{e \tau}{m} \vb{E}\)		&	Mobility	\\
			Photoelectric effect	&	\(	E = \hbar \nu\)	&	Planck's constant	\\
			\bottomrule
		\end{tabular}
	\end{table}
}
\end{frame}
%
\begin{frame}{\(y=mx\) and \(y= m /x\) relations contd...}
	%
		%
	\begin{itemize}
		\item Sometimes, physical quantities are also inversely related i.e. \(y \propto \tfrac{1}{x}\).
		\item We also say ``y is inversely proportional to x''.
	\end{itemize}
	%
	{
		%\small
		\begin{table}
			\centering
			\begin{tabular}{l c l}
				\toprule
				Law 	&	Equation 	&	Const. of prop.	\\
				\midrule
				Light	&	\(\nu = c\frac{1}{\lambda}\)	&	velocity	\\
				\bottomrule
			\end{tabular}
		\end{table}
	}
	%
\end{frame}
\begin{frame}{Mobility -- \(\mu\)}
	%
	\begin{definition}
		\textbf{Mobility} \(\mu\) is defined as the drift velocity per unit electric field
		\[
		\mu = \frac{v_d}{E}.
		\]
	\end{definition}
	%
	\begin{itemize}
		\item It is a material property as
		\[
		\mu = \frac{e\tau}{m} \qquad \qquad [\because \vb{v_d} = -\frac{e\tau}{m}\vb{E}]
		\]
		\item Its SI unit is
		\[
		[\mu] = \frac{[v_d]}{[E]} = \frac{[v_d \ell]}{[V]} = \text{\si{\metre\squared\per\volt\per\second}}
		\]
	\end{itemize}
	%
\end{frame}
%
\begin{frame}{Conductivity \(\sigma\)}
	%
	\begin{minipage}{0.28\textwidth}
		\begin{figure}
			\centering
			\includegraphics[width=\textwidth]{metals_conductivity_expression}
		\end{figure}
	\end{minipage}
	%
	\begin{minipage}{0.7\textwidth}
		In a time duration \(\Delta t\), the number of electrons \(N\) in a volume of \(V = v_d \Delta t A\) will pass through the cross section of area \(A\).
	\end{minipage}
	%
	\begin{small}
		\begin{derivation}
			%
			\begin{minipage}[t]{0.3\textwidth}
				\begin{enumerate}
					\item Current \(I\):
					\item Current density \(j\):
					\item Conductivity \(\sigma\):
				\end{enumerate}
			\end{minipage}
			\begin{minipage}[t]{0.68\textwidth}
				\vspace*{0pt} % forces content to start at top
				\begin{align}
					I &= -\frac{Ne}{\Delta t} =  -\frac{n\cdot V}{\Delta t} = - nev_d  A	\\
					%
					j &= \frac{I}{A} = -ne v_d = \dfrac{ne^2\tau}{m}E	\\
					%
					\sigma &= \dfrac{ne^2\tau}{m} = ne\big(\dfrac{e\tau}{m}\big) = ne\mu
				\end{align}
				
			\end{minipage}
		\end{derivation}
	\end{small}
	%
	\begin{learningobjectives}
		We have derived the expression for conductivity.
	\end{learningobjectives}
\end{frame}
%
\begin{frame}{Relaxation Time -- \(\tau\)}
	%
	\begin{block}{Assumption}
		Relaxation time \(\tau\) is assumed to be a \textbf{constant}.
	\end{block}
	\begin{guesstimate}{Relaxation time in metals}
		Estimate the electron relaxation time in metals.
		\begin{enumerate}\itemsep0pt\parskip0pt \parsep0pt
			\item Take conductivity \(\sigma \sim\)~\SI{e7}{\siemens\per\metre}, number density \(n \sim\)~\SI{e28}{\per\metre\cubed}
			\item Take electron mass  \(m\)~=\SI{9.1e-31}{\kilo\gram}, charge \(e\)~=\SI{1.609e-19}{\coulomb}.
			\item Deduce relaxation time \( \tau \sim\)~ \framebox{\SI{e-14}{\second} or \SI{10}{\femto\second}}.
		\end{enumerate}
	\end{guesstimate}
	\begin{itemize}
		\item Larger $\tau$ → fewer collisions → higher conductivity.
	\end{itemize}
	%
	\begin{trends}
		\begin{enumerate}\itemsep0pt\parskip0pt \parsep0pt
			\item As temperature increases, \(\tau\) decreases. \hfill Why?
			\item As impurity concentration increases, \(\tau\) decreases. \hfill Why?
		\end{enumerate}
	\end{trends}
\end{frame}
%
\begin{frame}{Resistivity -- \(\rho\)}
	\begin{definition}
		Resistivity \(\rho\) is defined as the inverse of conductivity
		\[
		\rho = \frac{1}{\sigma}
		\]
		Its SI unit is \si{\ohm\metre}.
	\end{definition}
	\begin{itemize}
		\item Resistivity $\rho = \dfrac{1}{\sigma} = \dfrac{m}{ne^2\tau}$.
		\item Its SI unit is \si{\ohm\metre}.
	\end{itemize}
	%
	\begin{trends}
		\begin{enumerate}\itemsep0pt\parskip0pt \parsep0pt
			\item As number density increases, $\rho$ decreases. \hfill Why?
			\item As temperature increases, $\rho$ increases. \hfill Why?
		\end{enumerate}
	\end{trends}
\end{frame}
%
\begin{frame}{Summary of physical quantities}
	\begin{table}
		\centering
		\begin{tabular}{l c c c}
			\toprule
			Physical quantity	&	Symbol	& Units	& Material Property	\\
			\midrule
			Resistance	&	\(R\)	&	\si{\ohm}	&	No	\\
			Conductivity	& \(\sigma\)	&	\si{\per\ohm\per\metre} or \si{\siemens\per\metre}	&	Yes	\\
			Mobility	&	\(\mu\)	&	\si{\metre\squared\per\volt\per\second}	&	Yes 	\\
			Relaxation Time	&	\(\tau\)	&	\si{\second}	& Yes	\\
			Resistivity	&	\(\rho\)	&	\si{\ohm\metre}	&	Yes	\\
			\bottomrule
		\end{tabular}
	\end{table}
\end{frame}
%
\begin{frame}{Theory vs. experiment}
	%
	\begin{block}{\textcolor{green}{Merit: Conductivity order of magnitude}}
		\begin{itemize}
			\item For metals, classical free electron theory (CEFT) predicts order of magnitude of conductivity correctly.  
			\[\sigma_{\mathrm{experiment}} \sim \sigma_{\mathrm{theory}} \sim \SI{e7}{\per\ohm\per\metre} \]
		\end{itemize}
	\end{block}
	%
	%
	\begin{block}{\textcolor{red}{Demerit: Conductivity vs valency}}
		\begin{itemize}
			\item Conductivity is proportional to number density. Number density is proportional to valency. So theory predicts conductivity increases with valency. But experiment shows monovalent metals (\ce{Na}, \ce{K}) have higher conductivity than divalent metals (\ce{Mg}, \ce{Ca}).
		\end{itemize}
	\end{block}
\end{frame}
%
\begin{frame}{Theory vs Experiment contd...}
	\begin{block}{\textcolor{red}{Demerit: Hall coefficient sign}}
		\begin{itemize}
			\item In metals, there are only electrons as charge carriers.
			\item It is expected that the \textbf{sign of Hall coefficient is negative}.
			\item However, experiments show metals like \ce{Zn} show positive Hall coefficient. CEFT cannot explain this \textbf{anomaly}.
		\end{itemize}
	\end{block}
	%
	\begin{block}{\textcolor{red}{Under-estimation of mean free path}}
		\begin{itemize}
			\item The mean free path is considered a constant in CEFT and is taken as the interatomic distance that is approximately \SI{5}{\angstrom}.
			\item However, experiment shows mean free path to be of the order \SI{50}{\angstrom}.
			\item Thus, CEFT underestimates the mean free path by one order of magnitude.
		\end{itemize}
	\end{block}
\end{frame}
\begin{frame}{\textcolor{red}{Demerit: Classification of materials}}
		%
		\begin{figure}
			\centering
			\includegraphics[width=0.7\textheight]{metals_semiconductors_resistivity}
		\end{figure}
		%
		\begin{itemize}
			\item CEFT \textbf{does not} explain the classification of materials into conductors, semi-conductors and insulators.
			\item For example, the resistivity of metals varies linearly
			\(
			\rho_\mathrm{conductor} \propto T
			\)
			but resistivity of semiconductors varies exponentially
			\(
			\rho_\mathrm{semi-conductor} \propto \exp(\frac{1}{T}).
			\)
		\end{itemize}
\end{frame}
%
\begin{frame}{\textcolor{red}{CFET Demerit: Resistivity vs temperature}}
	%
	\begin{figure}
		\centering
		\includegraphics[width=0.8\textheight]{metals_resistivity_theory_vs_experiment}
	\end{figure}
	%
	\begin{itemize}
		\item Wrong prediction of temperature dependence of conductivity. Theory predicts
		\(
		\sigma(T) \propto 1/\sqrt{T},
		\) 
		but experimental trend is
		\(
		\sigma(T) \propto T^{-1}.
		\)
	\end{itemize}
\end{frame}
%
\begin{frame}{\textcolor{red}{Demerit: Specific heat capacity}}
\begin{footnotesize}
		\begin{itemize}
		\item Metals heat ``quickly''. When we heat a metal, the heat energy is transferred to electrons. The quickness is measured in terms of heat capacity.
		\item Specific heat capacity is defined as the rate of change of energy per unit temperature per unit mole.
		\item Energy of electron at temperature \(T\) is
		\(
		E_\mathrm{el} = \frac{3}{2} k_B T.
		\)
		\item Energy of a mole of electrons is
		\[
		E = E_\mathrm{el} N_A =  \frac{3}{2} k_B N_A T  = \frac{3}{2} R T. \qquad \left[R = \SI{8.314}{\joule\per\mole\per\kelvin} \right]
		\]
		\item Therefore, specific heat capacity is
		\[
		C_\mathrm{theory} = \frac{\mathrm{d} E}{\mathrm{d} T} = \frac{3}{2} R \sim \SI{12}{\joule\per\mole\per\kelvin}.
		\]
		\item But experimental values of specific heat capacity \(C_\mathrm{experiment}\) are \textbf{smaller by two orders} of magnitude.
	\end{itemize}
\end{footnotesize}
%
\begin{keyinsight}
	Mercury thermometer is a good instrument to measure body temperature.
\end{keyinsight}
%
\end{frame}
%
%
\begin{frame}{Theory vs experiment contd...}
	\begin{block}{Reason}
		\begin{itemize}
			\item  Does not account for quantum effects (Fermi–Dirac statistics).
			\item Improved later by Sommerfeld's free electron theory (quantum).
		\end{itemize}
	\end{block}
\end{frame}
\begin{frame}{Summary of lecture}
	\begin{itemize}
		\item Derived $\sigma = \dfrac{ne^2\tau}{m}$.
		\item Connected microscopic electron properties to macroscopic Ohm’s law.
		\item Classical free electron theory has qualitative and quantitative agreement with experiment values of conductivity.
		\item However, the theory has drawbacks. Some of them are
		 \begin{itemize}
		 	\item Incorrect prediction for conductivity vs valency.
		 	\item Cannot explain anomalous sign of Hall coefficient in some metals.
		 	\item Underestimation of mean free path.
		 	\item Cannot explain classification of materials into conductors. semi-conductors and insulators.
		 	\item Wrong prediction of conductivity vs temperature.
		 	\item Overestimation of heat capacity.
		 \end{itemize} 
	\end{itemize}
\end{frame}